\chapter{Theoretical Background}
\label{ch:background}

The development of the Mesh-LZ codec is grounded in several well-established fields of information theory, statistical modeling, and digital signal processing. To fully contextualize the specific architectural choices made in Mesh-LZ, specifically its departure from traditional frequency-domain transforms in favor of spatial parsing, it is necessary to rigorously review the theoretical underpinnings of dictionary-based compression, modern entropy coding paradigms, and the mathematical properties of color decorrelation.

\section{Information Theory and Entropy}

The fundamental goal of any data compression system is to represent a given sequence of symbols using the theoretical minimum number of bits. In 1948, Claude Shannon established the foundation of information theory by defining the concept of Information Entropy, denoted as $H(X)$. For a discrete random variable $X$ with possible outcomes $x_i \in \{x_1, \dots, x_n\}$, each occurring with probability $P(x_i)$, the Shannon entropy is defined as:

\begin{equation}
    H(X) = - \sum_{i=1}^{n} P(x_i) \log_2 P(x_i)
\end{equation}

The entropy $H(X)$ represents the lower bound on the average number of bits required to encode the symbols generated by the source $X$. A fundamental property of this equation is that symbols with a high probability of occurrence convey less "surprise" or information, and therefore should be encoded using fewer bits. Conversely, rare symbols should be encoded using more bits.

\section{Entropy Coding Paradigms}

Entropy coders are the final stage of almost all compression pipelines. They take the raw symbols (or residuals) generated by the preceding modeling stages and convert them into a dense bitstream that approaches the Shannon entropy limit.

\subsection{Huffman Coding}
Developed by David A. Huffman in 1952, Huffman coding is a greedy algorithm that constructs an optimal prefix-free binary tree based on the probability of each symbol. The most frequent symbols are assigned the shortest binary codes (e.g., a single bit `0`), while less frequent symbols receive progressively longer codes.

While computationally inexpensive and easy to implement, Huffman coding suffers from a severe structural limitation: every symbol must be represented by an integer number of bits. This means the shortest possible code for any symbol is 1 bit. If a symbol is highly skewed and appears with a probability of $P(x) = 0.99$, its ideal theoretical bit-cost according to Shannon is $-\log_2(0.99) \approx 0.014$ bits. Huffman coding is forced to use 1 full bit, resulting in significant inefficiency (wasting nearly 0.986 bits per symbol) for highly skewed probability distributions.

\subsection{Arithmetic Coding}
Arithmetic Coding (AC) resolves the integer-bit limitation of Huffman coding. Instead of assigning a discrete binary code to each symbol, AC encodes an entire sequence of symbols into a single fractional number $F$ where $0 \le F < 1$. 

The algorithm operates by recursively subdividing the interval $[0, 1)$ according to the probability distribution of the symbols. As each symbol is processed, the current interval is narrowed to a sub-interval proportional to the probability of that symbol. When the entire sequence is processed, any binary fraction that falls within the final, microscopic interval uniquely identifies the original sequence.

Because AC can represent a symbol using fractional bits, it can reach the theoretical Shannon entropy limit almost perfectly. However, the continuous interval subdivisions require expensive fractional multiplications and divisions. Furthermore, the state is heavily serialized; decoding symbol $x_i$ requires the interval to have been fully updated by symbol $x_{i-1}$, severely restricting decoding speed on modern hardware.

\subsection{Asymmetric Numeral Systems (ANS)}
In 2009, Jarek Duda introduced Asymmetric Numeral Systems (ANS) \cite{duda2009asymmetric}, which revolutionized modern data compression (now utilized in Apple's LZFSE, Facebook's Zstandard, and JPEG XL). ANS bridges the performance gap, providing the exact compression ratios of Arithmetic Coding while operating at speeds comparable to or exceeding Huffman coding.

ANS operates by tracking a single natural number state variable $x$. To encode a symbol $s$ with a frequency $f_s$ within a total frequency scale $M$ (where $M = \sum f_s$), the state is updated such that the new state $x'$ mathematically encapsulates both the previous state $x$ and the new symbol $s$.

The encoding function for the Range variant of ANS (rANS) is defined mathematically as:
\begin{equation}
    x' = C(x, s) = \left\lfloor \frac{x}{f_s} \right\rfloor \times M + (x \pmod{f_s}) + C_s
\end{equation}
where $C_s$ is the cumulative frequency of all symbols strictly preceding $s$ in the alphabet: $C_s = \sum_{i < s} f_i$.

The corresponding decoding function extracting symbol $s$ and the previous state $x$ from the current state $x'$ is highly efficient:
\begin{equation}
    s = D(x' \pmod{M})
\end{equation}
\begin{equation}
    x = \left\lfloor \frac{x'}{M} \right\rfloor \times f_s + (x' \pmod{M}) - C_s
\end{equation}

Because the decoding process can be heavily optimized using hardware integer arithmetic, bitwise shifts (if $M$ is chosen as a power of 2), and lookup tables for $D()$, rANS enables extremely aggressive decompression. Mesh-LZ utilizes an Interleaved rANS implementation to maximize CPU pipeline saturation, which will be detailed in Chapter \ref{ch:algorithms}.

\section{Dictionary-Based Compression (Lempel-Ziv)}

The Lempel-Ziv family of algorithms (specifically LZ77 \cite{ziv1977universal}) operates on the principle of a sliding window dictionary. Rather than attempting to model the statistical probability of individual symbols like an entropy coder, LZ77 achieves compression by replacing repeated occurrences of data strings with references (pairs of \texttt{[offset, length]}) pointing to a previous occurrence of that string within a defined search window.

In standard text or binary data compression, the data is inherently one-dimensional. However, images are strictly two-dimensional topological grids. Spatial correlation in natural imagery occurs equally in horizontal, vertical, and diagonal directions. Traditional LZ77 applied naively to the raw byte stream of an image (e.g., raster scanning an uncompressed bitmap) fails to capture vertical correlation, as adjacent pixels in the Y-axis are physically separated by the entire row width (stride) in linear memory. 

PNG mitigates this issue by applying a set of predictive filters to decorrelate the 2D grid \textit{before} passing the residuals to the 1D LZ77 DEFLATE engine. Mesh-LZ proposes a fundamentally different methodology: rather than mathematically filtering the 2D grid, Mesh-LZ maps the 2D topology into a highly localized 1D path using space-filling curves and localized block stencils, thereby allowing the LZ dictionary to natively capture spatial repetitions without destructive predictive filtering.

\section{Color Space Decorrelation: The YCoCg-R Transform}

Natural images encoded in the standard RGB (Red, Green, Blue) color space exhibit massive inter-channel correlation. In most photographic imagery, a change in the intensity of a light source affects all three channels simultaneously. Encoding raw RGB values independently is highly inefficient.

To break this spectral redundancy, codecs typically transform the RGB data into a luma-chroma space, most commonly YCbCr (used heavily in JPEG and MPEG video). The YCbCr transform separates the brightness (Luma, Y) from the color differences (Chroma, Cb and Cr). However, the standard YCbCr transform requires floating-point matrix multiplications, which introduces rounding errors and prevents perfect mathematical reversibility, rendering it entirely unsuitable for strictly lossless compression modes.

Mesh-LZ employs the \textbf{YCoCg-R} (Luminance, Orange Chrominance, Green Chrominance - Reversible) color space, introduced by Malvar et al. \cite{malvar2003lifting}. Unlike YCbCr, the YCoCg-R transform can be computed entirely using computationally cheap integer additions, subtractions, and bitwise shifts. 

The forward YCoCg-R transform utilized in Mesh-LZ is defined mathematically as:
\begin{align}
    Co &= R - B \\
    t &= B + \left\lfloor \frac{Co}{2} \right\rfloor \\
    Cg &= G - t \\
    Y &= t + \left\lfloor \frac{Cg}{2} \right\rfloor
\end{align}

This transform is perfectly reversible, mapping the 24-bit RGB space precisely to the 24-bit YCoCg space (with one extra bit of dynamic range temporarily required for the chroma channels, handled via 16-bit signed integers internally). The inverse transform applied during decompression is:
\begin{align}
    t &= Y - \left\lfloor \frac{Cg}{2} \right\rfloor \\
    G &= Cg + t \\
    B &= t - \left\lfloor \frac{Co}{2} \right\rfloor \\
    R &= B + Co
\end{align}

This precise, bit-exact decorrelation allows Mesh-LZ to encode the Luma ($Y$) channel with high fidelity to preserve structural details, while the Chroma channels ($Co$, $Cg$) can optionally be aggressively quantized or spatially subsampled without severely impacting human psycho-visual perception.
